% ============================================================
% cita-iso690-html.tex — ISO 690 EN PANTALLA (SC-57)
% ------------------------------------------------------------
% LO CARGA SOLO EL DOCUMENTO MÍNIMO DE LA PASADA DE tex4ht,
% DESPUÉS DE cita-iso690-config.tex. CON ESTE ARCHIVO, HTML Y
% EPUB TOMAN EL TEXTO DE LAS CITAS Y DE LAS REFERENCIAS DE
% biblatex-iso690 (m_GenerarSalidasLibro.PrepararCitasPantalla).
%
% LO ÚNICO QUE CAMBIA RESPECTO DEL PDF: LOS NOMBRES CONSERVAN SU
% CAJA. EN EL PDF VAN \textsc{\MakeLowercase{...}}: EL TEXTO EN
% MINÚSCULA Y LA FUENTE LO PONE EN VERSALITAS. EN PANTALLA VAN
% \textsc{...} CON EL TEXTO COMO ESTÁ EN LA BASE, «Abdala, Omar»,
% Y LAS VERSALITAS LAS PONE EL CSS (all-small-caps). UN LECTOR
% QUE NO LAS RESPETE MUESTRA EL NOMBRE COMO ESTÁ EN LA BASE Y NO
% «abdala, omar». ES LA MISMA DECISIÓN QUE VERONA (SC-56).
%
% LAS DEFINICIONES DEL PDF VAN EN \familynameformat Y EN UN
% \AtBeginBibliography; ESTAS SE AGREGAN DESPUÉS Y LAS PISAN.
%
% backref=false YA LO CARGA cita-iso690-config.tex. LAS NOTAS:
% EN EL CAMINO DE ISO 690 \iffootnote SOLO LO USA \smartcite,
% QUE LA APLICACIÓN NO EMITE (authoryear.cbx, biblatex 3.19).
% ============================================================
\renewcommand{\familynameformat}[1]{\textsc{#1}}

\AtBeginBibliography{%
	\renewcommand*{\mkbibnamegiven}[1]{\textsc{#1}}%
	\renewcommand*{\mkbibnameprefix}[1]{\textsc{#1}}%
	\renewcommand*{\mkbibnamesuffix}[1]{\textsc{#1}}%
}

% ------------------------------------------------------------
% \gbImprimirBibliografia — LA EMITE docbook-to-latex-citas.xsl
% ------------------------------------------------------------
% SIN TÍTULO: EL DE LA BIBLIOGRAFÍA LO PONE LA HOJA DE HTML O EPUB
% ------------------------------------------------------------
\newcommand*{\gbImprimirBibliografia}{%
  \printbibliography[heading=none]}
